\section*{第 \ref{ch:b1:radiometric-methods} 章——放射性测年方法}
% values: calc/ch25.py

\begin{solution}{exo:b1:radiometric-methods:1}
(a)~$D^*/P = 2^n - 1 = 3$ 得 $n = 2$：即两个半衰期。(b)~铷-87：
$e^{\lambda t} - 1 = e^{\num{1.3972e-11} \times \num{3.10e8}} - 1 = 0.004341$;
铀-235：$e^{\num{9.8485e-10} \times \num{3.10e8}} - 1 = 0.3570$。
铀的比值约大 80 倍：同样的测量绝对误差，
在其中所占的比例小得多。(c)~$2^{-n} < 0.01$，当 $n > \log_2 100 = 6.64$ 时：
约七个半衰期之后。
\end{solution}

\begin{solution}{exo:b1:radiometric-methods:2}
(a)~对木炭做放射性碳测年。(b)~对锆石做铀–铅法，或
对透\omterm{def:b1:rock-forming-minerals:groups}{长石}做氩–氩法（钾–氩法），透\omterm{def:b1:rock-forming-minerals:groups}{长石}是一种富钾\omterm{def:b1:rock-forming-minerals:mineral}{矿物}。
(c)~锆石适于铀–铅法，它能经受漫长的历史；铷–锶法
用于全岩也可行；钾–氩法则有使年龄偏低的风险，原因是
后期加热时氩丢失。(d)~对贝壳碳酸盐用放射性碳法，
接近该方法的极限（约 \num{70000} 年）。(e)~没有直接方法：
其年龄由上下两侧的定年火山岩限定，或
将其化石与已测年的剖面对比而定（见\cref{ch:b1:stratigraphic-scale}）。
\end{solution}

\begin{solution}{exo:b1:radiometric-methods:3}
(a)~$t = -8033 \ln 0.25 = 8033 \times 1.3863 = \num{11136}$ 年 BP。
(b)~平均寿命为 $5700/\ln 2 = \num{8223}$ 年：$t = 8223 \times 1.3863 =
\num{11400}$ 年。(c)~两种计算都假定，当时生物碳中碳-14 的比例
与标准值相同；事实并非如此。日历年龄由
\omterm{def:b1:radiometric-methods:radiocarbon}{标准曲线}（如 IntCal20）读出，以 cal BP 年的区间表示。
\end{solution}

\begin{solution}{exo:b1:radiometric-methods:4}
(a)~每个“时钟”在\omterm{def:b1:rock-forming-minerals:mineral}{矿物}冷却至封闭温度以下时启动：
锆石（高于 \qty{750}{\celsius}）最先闭合，\omterm{def:b1:rock-forming-minerals:groups}{角闪石}（约
\qty{530}{\celsius}）次之，黑\omterm{def:b1:rock-forming-minerals:groups}{云母}（约 \qty{280}{\celsius}）最后，因此年龄
按此顺序递减，反映花岗岩的冷却过程。(b)~$(530 - 280)/(410 - 398) =
250/12 = \qty{21}{\celsius}$ 每百万年。(c)~锆石在高于
\qty{750}{\celsius}时闭合，因此花岗岩至少冷却了 $750 - 530 =
\qty{220}{\celsius}$，历时 \qty{10}{\Ma}，即每百万年至少
\qty{22}{\celsius}，这是一个下限。
\end{solution}

\begin{solution}{exo:b1:radiometric-methods:5}
(a)~$\bar x = 7.5$，$\bar y = 0.809525$，$S_{xx} = 81.0$，$S_{xy} = 1.13965$；
$m = 1.13965/81.0 = 0.0140698$, $b = 0.809525 - 7.5 \times 0.0140698 = 0.70400$.
(b)~$t = \ln(1.0140698)/\num{1.3972e-11} = \qty{1000.0}{\Ma}$；\omterm{def:b1:radiometric-methods:isochron}{初始比值}
$0.7040$。(c)~残差为 $+1.6$、$-2.5$、$+0.7$ 和 $+0.2$（单位为
\num{e-4}): $s = \sqrt{\num{9.34e-8}/2} = 0.00022$, $\sigma_m = 0.00022/\sqrt{81}
= \num{2.40e-5}$ 和 $\sigma_t = \num{2.40e-5}/(\num{1.3972e-11} \times 1.0141) =
\qty{1.7}{\Ma}$：\qty[uncertainty-mode=separate]{1000.0(17)}{\Ma}。(d)~$0.7040$
接近\omterm{def:b1:layered-earth:layers}{地幔}值 $0.7025$，且远低于\omterm{def:b1:layered-earth:layers}{地壳}值 $0.710$：源区
主要位于\omterm{def:b1:layered-earth:layers}{地幔}中。
\end{solution}

\begin{solution}{exo:b1:radiometric-methods:6}
(a)~$0.100/39.098 \times \num{1.17e-4} = \qty{2.992e-7}{mol}$，即每
克钾-40 的物质的量。(b)~${}^{40}\mathrm{Ar}^*/{}^{40}\mathrm{K} = \num{1.169e-9}/\num{2.992e-7}
= 0.0039065$; $\lambda/\lambda_e = 5.543/0.581 = 9.5404$; $t = \ln(1 + 9.5404
\times 0.0039065)/\num{5.543e-10} = \ln(1.037269)/\num{5.543e-10} =
\qty{66.01}{\Ma}$。(c)~若放射成因氩将为 $0.995 \times 1.169$
纳摩尔，得 \qty{65.69}{\Ma}：实测年龄将偏大 \qty{0.32}{\Ma}
偏大，是 \num{32000} 年时限的十倍，该时限内撞击与
白垩纪末的灭绝已被证明同时发生。过剩氩主要是年轻岩石和全岩的问题；
透\omterm{def:b1:rock-forming-minerals:groups}{长石}之类的\omterm{def:b1:rock-forming-minerals:mineral}{矿物}在结晶时不应吸收氩，
氩–氩年龄谱会将其揭示出来。
\end{solution}

\begin{solution}{exo:b1:radiometric-methods:7}
(a)~约 \num{2000} cal BP：曲线在该处经过 \num{2070} 年 BP。
(b)~约 \num{2495} 至 \num{2699} cal BP，相差两百年：曲线在
\num{2500} 至 \num{2700} cal BP 之间近乎平直，因此该区间内每个日历年龄
对应的\omterm{def:b1:radiometric-methods:radiocarbon}{放射性碳年龄}几乎相同，更精确的
测量也无法缩小这一范围。(c)~在 \num{2000} cal BP 处，\omterm{def:b1:radiometric-methods:radiocarbon}{放射性碳
年龄}比日历年龄\emph{大} \num{70} 年；在 \num{8000} cal BP 处
\omterm{def:b1:radiometric-methods:radiocarbon}{放射性碳年龄}小 \num{822} 年（\qty{10}{\percent}）；在 \num{12000} cal BP 处小 \num{1694}
年（\qty{14}{\percent}）。差值与年龄并不成正比，
甚至符号也会改变。它不可能仅由半衰期造成，否则
将给出恒定的 \qty{2.4}{\percent}。
\end{solution}

\begin{solution}{exo:b1:radiometric-methods:8}
(a)~下方的火山灰层（\qty{94.50}{\Ma}）给出最大年龄，上方的熔岩流
（\qty{93.80}{\Ma}）给出最小年龄：界线限定在两者之间。岩墙也
给出最小年龄，但约束较弱，比熔岩更年轻。碎屑锆石仅对砂岩
给出最大年龄，且弱于火山灰层，后者位于
砂岩之下。(b)~$f = (12 - 4)/(20 - 4) = 0.5$；$t = 94.50 - 0.5 \times 0.70 =
\qty{94.15}{\Ma}$; $\sigma_t = \sqrt{0.25 \times 0.10^2 + 0.25 \times 0.12^2} =
\qty{0.08}{\Ma}$：\qty[uncertainty-mode=separate]{94.15(8)}{\Ma}。(c)~其
颗粒来自至少 \qty{95.3}{\Ma} 的岩石，该岩石在源区
某处遭受侵蚀；这说明砂岩的年龄小于 \qty{95.3}{\Ma}，这一点
下方的火山灰层已更严格地证明。
\end{solution}

\begin{solution}{exo:b1:radiometric-methods:9}
(a)~颗粒~A：$\ln(1.08401)/\num{1.55125e-10} = \qty{520.0}{\Ma}$ 和
$\ln(1.6688)/\num{9.8485e-10} = \qty{520.0}{\Ma}$。颗粒~B：\qty{368.4}{\Ma} 和
\qty{389.9}{\Ma}。(b)~颗粒~A 谐和。(c)~$0.05881/0.08401 = 0.7000$ 和
$0.4682/0.6688 = 0.7001$。年龄为 $\ln(1 + R)/\lambda$，并不正比于
$R$：对于较小的铀-238 比值 $\ln(1+R) \approx R$，年龄下降
近 \qty{30}{\percent}（至 \qty{368}{\Ma}）；对于较大的铀-235 比值
对数下降较少，年龄仅降至 \qty{390}{\Ma}。两者都偏小，
且二者互不一致：该颗粒不谐和。
\end{solution}

\begin{solution}{exo:b1:radiometric-methods:10}
(a)~令 $u = \lambda t$，$D^*/P = e^u - 1 = u + u^2/2 + \dots$，则 $(D^*/P)/\lambda
= t\,(1 + u/2 + \dots)$：近似高估年龄，相对
误差接近 $u/2$。(b)~铷-87 在 \qty{310}{\Ma}：$u = 0.00433$，误差
\qty{0.22}{\percent}（\qty{0.7}{\Ma}）；钾-40 在 \qty{296}{\Ma}：$u =
0.164$，误差 \qty{8.7}{\percent}（\qty{26}{\Ma}；$u/2$ 给出 \qty{8.2}{\percent}）；
铀-238 在 \qty{4.4}{\Ga}：$u = 0.683$，误差 \qty{43}{\percent}——近似
乃至其误差估计（$u/2 = \qty{34}{\percent}$）均失效。
(c)~$R = 0.004341$，$\sigma_R = \num{4.34e-6}$；由
\cref{prop:b1:radiometric-methods:precision}可得 $\sigma_t =
\num{4.34e-6}/(\num{1.3972e-11} \times 1.0043) = \qty{0.31}{\Ma}$.
\end{solution}

\begin{solution}{exo:b1:radiometric-methods:11}
(a)~设 ${}^{40}\mathrm{K}$ 为今天的钾-40。在距今 $t_2$ 的时刻
共有 ${}^{40}\mathrm{K}\,e^{\lambda t_2}$ 个原子，而
在 $t_1$ 与 $t_2$ 之间生成的氩，由\cref{prop:b1:radiometric-methods:kar}应用于
该区间，其值为 $(\lambda_e/\lambda)\,{}^{40}\mathrm{K}\,e^{\lambda
t_2}\bigl(e^{\lambda (t_1 - t_2)} - 1\bigr) = (\lambda_e/\lambda)\,
{}^{40}\mathrm{K}\,\bigl(e^{\lambda t_1} - e^{\lambda t_2}\bigr)$。其中留存的
比例为 $1 - f$。自 $t_2$ 以来，封闭\omterm{def:b1:rock-forming-minerals:mineral}{矿物}又生成了
$(\lambda_e/\lambda)\,{}^{40}\mathrm{K}\,(e^{\lambda t_2} - 1)$ 的氩。
将二者相加，再乘以 $\lambda/(\lambda_e\,{}^{40}\mathrm{K})$，即得该表达式。
(b)~表观年龄为 $\ln(1 + \text{表达式})/\lambda$，依次为
\qty{400.0}{\Ma}（$f = 0$）、\qty{342.3}{\Ma}（$0.3$）、\qty{282.7}{\Ma}（$0.6$）和
\qty{200.0}{\Ma}（$1$）。(c)~介于 \qty{200}{\Ma} 与
\qty{400}{\Ma} 之间的每一个年龄都可由某种部分丢失产生，单个数值无法
判断它记录的是该时刻的真实事件，还是混合结果；需要年龄谱，
或同一岩石中的其他计时体系。
\end{solution}

\begin{solution}{exo:b1:radiometric-methods:12}
(a)~年龄与所用平均寿命成正比：$10306 \times 8223/8033 =
\num{10550}$ 年。(b)~$10550 - 10306 = 244$ 年，约为
\num{1694} 的七分之一。(c)~若初始活度为 $k A_0$，则真实年龄为 $t = \tau
\ln(k A_0/A) = \num{10550} + \tau \ln k$；当 $t = \num{12000}$ 时，$\ln k =
1450/8223 = 0.176$，且 $k = 1.19$：碳-14 约多 \qty{19}{\percent}，
比标准样品中的含量高。
\end{solution}

\begin{solution}{pb:b1:radiometric-methods:1}
\begin{enumerate}
  \item 各点随 $x$ 稳步上升；它们偏离某直线的偏差至多为
        小数点后第四位上的两个单位，而 $0.04$ 是 $y$ 的变化范围：它们大致位于
        一条直线上，散布很小。
  \item $\bar x = 4.90333$, $\bar y = 0.729288$, $S_{xx} = 62.1253$, $S_{xy} =
        0.269335$; $m = 0.269335/62.1253 = 0.0043354$.
  \item $b = 0.729288 - 0.0043354 \times 4.90333 = 0.70803$：即
        花岗岩结晶时的 $^{87}\mathrm{Sr}/^{86}\mathrm{Sr}$，
        为其全部部分所共有。
  \item $t = \ln(1.0043354)/\num{1.3972e-11} = \qty{309.6}{\Ma}$.
  \item 残差依次为 $+9.1$、$-17.2$、$+6.9$、$-6.4$、$+15.1$ 和 $-7.5$（以
        \num{e-5} 为单位）；$s = 0.000137$，$\sigma_m = 0.000137/\sqrt{62.1253} =
        \num{1.74e-5}$, $\sigma_t = \num{1.74e-5}/(\num{1.3972e-11} \times 1.0043)
        = \qty{1.2}{\Ma}$：\qty[uncertainty-mode=separate]{309.6(12)}{\Ma}，且
        $\sigma_b = 0.00010$。
  \item $0.7080$ 更接近\omterm{def:b1:layered-earth:layers}{地壳}的 $0.710$，而不是\omterm{def:b1:layered-earth:layers}{地幔}的 $0.7025$。
        若为混合物，$(0.70803 - 0.7025)/(0.710 - 0.7025) = 0.74$：约
        四分之三的锶来自熔融的\omterm{def:b1:layered-earth:layers}{地壳}。
  \item $0.0740/39.098 \times \num{1.17e-4} = \qty{2.214e-7}{mol}$ 每克，
        乘以 $\num{6.022e23}$：\num{1.33e17} 个原子。
  \item $^{40}\mathrm{Ar}^*/^{40}\mathrm{K} = \num{4.139e-9}/\num{2.2144e-7} =
        0.018691$; $1 + 9.5404 \times 0.018691 = 1.17832$; $t =
        \ln(1.17832)/\num{5.543e-10} = \qty{296.0}{\Ma}$.
  \item $^{40}\mathrm{K} = \qty{3.292e-8}{mol.g^{-1}}$，比值 $0.019239$，
        $\ln(1.18355)/\num{5.543e-10} = \qty{304.0}{\Ma}$.
  \item $\ln(1.018691)/\num{5.543e-10} = \qty{33.4}{\Ma}$，年龄偏小了九倍，
        因为只有约十分之一的衰变产生氩。钙-40 是钙的常见同位素，
        在岩石中丰度很高：由衰变产生的少量
        无法与\omterm{def:b1:rock-forming-minerals:mineral}{矿物}形成之初便含有的钙相区分。
  \item 过剩氩使年龄偏移的量级约为百万年，这对年轻岩石
        很重要；此处影响远低于年龄的 \qty{1}{\percent}。
        此外，\omterm{def:b1:rock-forming-minerals:groups}{云母}与\omterm{def:b1:rock-forming-minerals:groups}{角闪石}晶体在
        形成时几乎不吸收氩。
  \item 或是花岗岩缓慢冷却，\omterm{def:b1:rock-forming-minerals:groups}{云母}与\omterm{def:b1:rock-forming-minerals:groups}{角闪石}在
        结晶之后才封闭（\omterm{def:b1:radiometric-methods:closure}{冷却年龄}），或是后来的加热驱出了它们的部分
        氩。在氩–氩年龄谱中，平坦的坪段支持\omterm{def:b1:radiometric-methods:closure}{冷却年龄}；
        无坪段的年龄谱则提示\omterm{def:b1:rock-forming-minerals:mineral}{矿物}丢失了氩。
        这些年龄的先后次序，即普通\omterm{def:b1:rock-forming-minerals:groups}{角闪石}先于黑\omterm{def:b1:rock-forming-minerals:groups}{云母}，
        与其封闭温度的预测一致，同样指向冷却。
  \item $\ln(1.049361)/\num{1.55125e-10} = \qty{310.6}{\Ma}$ 与
        $\ln(1.35784)/\num{9.8485e-10} = \qty{310.6}{\Ma}$：二者协和。
  \item 锆石 \qty{310.6}{\Ma}，封闭温度高于 \qty{750}{\celsius}：
        结晶时间；全岩\omterm{def:b1:radiometric-methods:isochron}{等时线}，
        \qty[uncertainty-mode=separate]{309.6(12)}{\Ma}，同样测定了它的年龄，
        并与之在不确定度范围内一致；普通\omterm{def:b1:rock-forming-minerals:groups}{角闪石} \qty{304.0}{\Ma}：冷却
        经过约 \qty{530}{\celsius}；黑\omterm{def:b1:rock-forming-minerals:groups}{云母} \qty{296.0}{\Ma}：冷却
        经过约 \qty{280}{\celsius}。图中曲线在
        \qty{310.6}{\Ma} 时高于 \qty{750}{\celsius}，
        至 \qty{296.0}{\Ma} 时降至 \qty{280}{\celsius}。
  \item $(530 - 280)/(304.02 - 296.03) = 250/7.99 = \qty{31}{\celsius}$ 每
        百万年。
  \item $(750 - 530)/(310.60 - 304.02) = 220/6.58 = \qty{33}{\celsius}$ 每
        百万年，至少为此值，因为锆石在高于
        \qty{750}{\celsius} 时即已封闭：早期的冷却至少同样快。
  \item 温度差 $250 \pm \sqrt{40^2 + 40^2} = 250 \pm
        57$，给出 \qtyrange{24}{38}{\celsius} 每百万年。
  \item $(550 - 350)/7.99 = \qty{25}{\celsius}$ 每百万年，落在
        第~17~题给出的范围内：不确定度来自封闭温度，而非
        年龄。
  \item $(280 - 10)/30 = \qty{9.0}{km}$.
  \item $\ln(1.039108)/\num{1.55125e-10} = \qty{247.30}{\Ma}$，由
        铀-235 比值得到相同结果：协和。
  \item 来自在\omterm{def:b1:cross-sections:unconformity}{不整合}面处被剥蚀的花岗岩：其年龄即其
        结晶年龄。它们仅说明砂岩比
        \qty{310}{\Ma} 年轻：这只是最大年龄，约束力远不如火山灰层。
  \item $(58 - 18)/(247.30 - 245.30) = \qty{20.0}{m}$ 每百万年；层底
        比下部火山灰层老 $18/20.0 = \qty{0.90}{\Ma}$：
        \qty{248.20}{\Ma}.
  \item 由 \qty{309.6}{\Ma}（\omterm{def:b1:radiometric-methods:isochron}{等时线}）或 \qty{310.6}{\Ma}（锆石）至
        \qty{248.2}{\Ma}：约 \qty{62}{\Ma}。
  \item \qty{9.0}{km} 在 $296.0 - 248.2 = \qty{47.8}{\Ma}$ 内：\qty{0.19}{km} 每
        百万年，即 \qty{0.19}{mm} 每年。这一结论假定，所测的
        黑\omterm{def:b1:rock-forming-minerals:groups}{云母}取自紧邻\omterm{def:b1:cross-sections:unconformity}{不整合}面之下，因而花岗岩与片岩中位于其上的全部岩石
        在沙沉积之前已被剥蚀，
        且\omterm{def:b1:heat-budget:geotherm}{地温线}保持恒定。
  \item 花岗岩以约 \qty{31}{\celsius} 每百万年的速率冷却，
        即在普通\omterm{def:b1:rock-forming-minerals:groups}{角闪石}与黑\omterm{def:b1:rock-forming-minerals:groups}{云母}封闭之间。
\end{enumerate}
\end{solution}
